% =============================================================================
% SECTION 3: KNOWLEDGE GRAPH DESIGN
% =============================================================================

\section{Knowledge Graph Design}
\label{sec:kg}

The knowledge graph (KG) serves as the backbone for \emph{grounded, repeatable decisions}~\cite{hogan2021kg_survey, zheng2022kg_construction}.
Unlike rule-based systems that embed logic in code, the KG externalizes knowledge so it can be inspected, updated, and audited independently of the reasoning engine.

\subsection{Theoretical Foundation}

Knowledge graphs represent information as a labeled, directed multi-graph $G = (V, E, L)$ where vertices $V$ represent entities, edges $E \subseteq V \times V$ represent relationships, and $L$ provides type labels~\cite{hogan2021kg_survey}.
This formalism supports:

\begin{itemize}
  \item \textbf{Explicit semantics}: Relationships are first-class citizens with types and properties
  \item \textbf{Schema flexibility}: New entity and relationship types can be added without migration
  \item \textbf{Path queries}: Multi-hop reasoning (e.g., ``which code rules constrain this header?'')
  \item \textbf{Provenance tracking}: Every fact can carry source attribution
\end{itemize}

Pan et al.~\cite{pan2024unifying_kg_llm} identify three paradigms for KG-LLM integration: KG-enhanced LLMs, LLM-augmented KGs, and synergized systems.
Our architecture follows the synergized approach, where agents query the KG for grounded facts and update it with validated inferences.

\subsection{Why Not Alternatives?}

We considered three alternative knowledge representations:

\textbf{Rule engines} (e.g., Drools, CLIPS) excel at forward/backward chaining but embed rules in code or DSL files.
Rules are difficult to inspect, version, and explain to end users.
When rules conflict, debugging requires developer intervention.

\textbf{Ontologies} (e.g., OWL/RDF) provide formal semantics and reasoning via description logics~\cite{pauwels2017kg_ifc}.
However, OWL reasoners struggle with scale, closed-world assumptions conflict with construction's open-world nature, and the formalism is inaccessible to domain experts.

\textbf{Relational databases} handle structured data well but require JOIN-heavy queries for relationship traversal.
A query like ``find all components constrained by IRC R602.3'' requires multiple joins across normalized tables; in a graph, it is a single-hop traversal.

The KG approach balances expressiveness, scalability, and maintainability for construction domain knowledge.

\subsection{Knowledge Categories}

The KG organizes construction knowledge into six categories:

\begin{itemize}
  \item \textbf{Component definitions}: studs, headers, plates, cripples, blocking, sheathing---with parameterized assemblies
  \item \textbf{Material catalog}: dimensions, grades, treatments, SKUs, stock lengths, pricing, lead times
  \item \textbf{Fasteners and connectors}: nails, screws, straps, hangers, anchors---with application rules and schedules
  \item \textbf{Building rules}: framing conventions, span tables, loading assumptions, nailing patterns
  \item \textbf{Codes and regulations}: jurisdiction, version, section citations, applicability conditions
  \item \textbf{Historical precedents}: past project decisions, substitutions, waste rates, outcome feedback
\end{itemize}

\subsection{Core Entities and Relationships}

Table~\ref{tab:kg_entities} summarizes the core entity types, and Table~\ref{tab:kg_relationships} shows key relationships.

\begin{table}[ht]
\centering
\caption{Core KG Entity Types}
\label{tab:kg_entities}
\begin{tabular}{ll}
\toprule
\textbf{Entity Type} & \textbf{Description} \\
\midrule
Project & Building project with metadata \\
PlanSheet & Individual plan page with scale \\
PlanFact & Extracted fact with confidence \\
AssemblyIntent & Inferred assembly (e.g., header) \\
Component & Physical component instance \\
StockItem & Purchasable material \\
CutPiece & Result of cutting operation \\
CodeRule & Building code requirement \\
FastenerRule & Fastening specification \\
Supplier & Material supplier \\
\bottomrule
\end{tabular}
\end{table}

\begin{table}[ht]
\centering
\caption{Core KG Relationship Types}
\label{tab:kg_relationships}
\begin{tabular}{lll}
\toprule
\textbf{Relationship} & \textbf{From} & \textbf{To} \\
\midrule
hasSheet & Project & PlanSheet \\
mentions & PlanSheet & PlanFact \\
instantiates & PlanFact & AssemblyIntent \\
composedOf & AssemblyIntent & Component \\
builtFrom & Component & StockItem \\
cutInto & StockItem & CutPiece \\
constrainedBy & Component & CodeRule \\
requiresFastener & Component & FastenerRule \\
suppliedBy & StockItem & Supplier \\
\bottomrule
\end{tabular}
\end{table}

\subsection{Example Query: Provenance Chain}

The following Cypher query retrieves the full provenance chain for a BOM line item, demonstrating the KG's traceability:

\begin{lstlisting}[language=json, basicstyle=\ttfamily\scriptsize]
MATCH (bom:BOMItem {id: $itemId})
      -[:derivedFrom]->(cut:CutPiece)
      -[:cutFrom]->(stock:StockItem)
      <-[:builtFrom]-(comp:Component)
      <-[:composedOf]-(asm:AssemblyIntent)
      <-[:instantiates]-(fact:PlanFact)
      -[:extractedFrom]->(sheet:PlanSheet)
OPTIONAL MATCH (comp)-[:constrainedBy]->(code:CodeRule)
RETURN bom, cut, stock, comp, asm, fact,
       sheet, collect(code) as codeRules
\end{lstlisting}

This single query traverses the entire chain from BOM output back to plan source, collecting all applicable code rules---a query that would require 6+ JOINs in a relational model.

\subsection{Versioning and Provenance}

Construction knowledge evolves: codes are updated, suppliers change SKUs, practices differ by region.
The KG supports this through:

\begin{itemize}
  \item \textbf{Temporal validity}: Effective dates and regions on CodeRule and MaterialSpec nodes
  \item \textbf{Source provenance}: Document reference, section number, URL/file hash on every fact
  \item \textbf{Decision audit}: Agent decisions are logged with reasoning traces linked to KG facts
\end{itemize}

Every generated BOM line item links back to: (1) which plan facts triggered it, (2) which inference rules were applied, and (3) which code references constrained it.

\subsection{Technology Choice: Neo4j}

We select Neo4j as the graph database platform based on:

\begin{itemize}
  \item \textbf{Native graph storage}: Index-free adjacency enables O(1) relationship traversal regardless of graph size~\cite{neo4j2023construction}
  \item \textbf{Cypher query language}: Declarative pattern matching for complex multi-hop queries
  \item \textbf{Vector index support}: Native vector similarity search enables GraphRAG patterns~\cite{edge2024graphrag, lewis2020rag}
  \item \textbf{APOC procedures}: Graph algorithms (centrality, community detection) for analytics
  \item \textbf{Python ecosystem}: Async driver compatible with FastAPI; LangChain integration available
\end{itemize}

Alternative graph databases (Amazon Neptune, TigerGraph) were considered but Neo4j's developer tooling, community support, and vector search capabilities made it the pragmatic choice for this domain.
